% Options for packages loaded elsewhere
\PassOptionsToPackage{unicode}{hyperref}
\PassOptionsToPackage{hyphens}{url}
\documentclass[
  ignorenonframetext,
]{beamer}
\newif\ifbibliography
\usepackage{pgfpages}
\setbeamertemplate{caption}[numbered]
\setbeamertemplate{caption label separator}{: }
\setbeamercolor{caption name}{fg=normal text.fg}
\beamertemplatenavigationsymbolsempty
% remove section numbering
\setbeamertemplate{part page}{
  \centering
  \begin{beamercolorbox}[sep=16pt,center]{part title}
    \usebeamerfont{part title}\insertpart\par
  \end{beamercolorbox}
}
\setbeamertemplate{section page}{
  \centering
  \begin{beamercolorbox}[sep=12pt,center]{section title}
    \usebeamerfont{section title}\insertsection\par
  \end{beamercolorbox}
}
\setbeamertemplate{subsection page}{
  \centering
  \begin{beamercolorbox}[sep=8pt,center]{subsection title}
    \usebeamerfont{subsection title}\insertsubsection\par
  \end{beamercolorbox}
}
% Prevent slide breaks in the middle of a paragraph
\widowpenalties 1 10000
\raggedbottom
\AtBeginPart{
  \frame{\partpage}
}
\AtBeginSection{
  \ifbibliography
  \else
    \frame{\sectionpage}
  \fi
}
\AtBeginSubsection{
  \frame{\subsectionpage}
}
\usepackage{iftex}
\ifPDFTeX
  \usepackage[T1]{fontenc}
  \usepackage[utf8]{inputenc}
  \usepackage{textcomp} % provide euro and other symbols
\else % if luatex or xetex
  \usepackage{unicode-math} % this also loads fontspec
  \defaultfontfeatures{Scale=MatchLowercase}
  \defaultfontfeatures[\rmfamily]{Ligatures=TeX,Scale=1}
\fi
\usepackage{lmodern}
\ifPDFTeX\else
  % xetex/luatex font selection
\fi
% Use upquote if available, for straight quotes in verbatim environments
\IfFileExists{upquote.sty}{\usepackage{upquote}}{}
\IfFileExists{microtype.sty}{% use microtype if available
  \usepackage[]{microtype}
  \UseMicrotypeSet[protrusion]{basicmath} % disable protrusion for tt fonts
}{}
\makeatletter
\@ifundefined{KOMAClassName}{% if non-KOMA class
  \IfFileExists{parskip.sty}{%
    \usepackage{parskip}
  }{% else
    \setlength{\parindent}{0pt}
    \setlength{\parskip}{6pt plus 2pt minus 1pt}}
}{% if KOMA class
  \KOMAoptions{parskip=half}}
\makeatother
\setlength{\emergencystretch}{3em} % prevent overfull lines
\providecommand{\tightlist}{%
  \setlength{\itemsep}{0pt}\setlength{\parskip}{0pt}}
\usepackage{bookmark}
\IfFileExists{xurl.sty}{\usepackage{xurl}}{} % add URL line breaks if available
\urlstyle{same}
\hypersetup{
  hidelinks,
  pdfcreator={LaTeX via pandoc}}

\author{\texorpdfstring{}{}}
\date{}

\begin{document}

\begin{frame}[fragile]{Problem statement}
\protect\phantomsection\label{problem-statement}
Fix an integer \(Q \geq 1\). For a real \(x\), define \begin{equation}
\label{eq:delta}
\delta_Q(x) = \min_{p \in \mathbb{Z},\, 1 \leq q \leq Q} \left| x - \frac{p}{q} \right|.
\end{equation}

The minimum exists because for each \(q\) the optimal numerator is one
of three integers around \(qx\) (convexity of \(p \mapsto |x-p/q|\)
along \(\mathbb{Z}\)), so the search over \(p\) at fixed \(q\) reduces
to checking \(p_0-1,p_0,p_0+1\) with \(p_0=\mathrm{round}(qx)\).
Enumerating \(q=1,\ldots,Q\) yields an \(O(Q)\) exact procedure in
floating-point arithmetic at the stated grid.

\begin{block}{Circle formulation}
\protect\phantomsection\label{circle-formulation}
For analysis on \(\mathbb{R}/\mathbb{Z}\) one often studies the
fractional part \(\{x\}\) and replaces \(x\) by \(\{x\}\) in
\eqref{eq:delta} when comparing to laws supported on \([0,1)\). The
codebase exposes both the raw distance on \(\mathbb{R}\) and the
fractional variant when \texttt{experiment.compare\_mod1} is true in
\texttt{manuscript/config.yaml}.
\end{block}

\begin{block}{Statistical question}
\protect\phantomsection\label{statistical-question}
\textbf{Research question (statistical).} Conditional on a chosen
reference law for \(X\) (here: a mixture of \(\mathrm{Unif}(0,1)\),
fractional parts \(\{\sqrt{k}\}\) with \(k\) sampled from a fixed
square-free list, and optionally \(\mathrm{Beta}(a,b)\) with parameters
from config), where do deterministic constants \(x^\star\) sit in the
empirical distribution of \(\delta_Q(X)\)? Equivalently: what empirical
percentile does \(\delta_Q(x^\star)\) attain relative to Monte Carlo
samples at the same \(Q\)? The same question applies to \(\mu_Q\) when
reporting denominator-weighted quality (\eqref{eq:lagrange} in
\texttt{03e\_q\_squared\_quality.md}).
\end{block}

\begin{block}{Rational control}
\protect\phantomsection\label{rational-control}
Any \(x \in \mathbb{Q}\) yields \(\delta_Q(x)=0\) once \(Q\) is at least
the reduced denominator of \(x\), and \(\mu_Q(x)=0\) simultaneously. The
registry includes \texttt{one\_sixth} as a \textbf{structural sanity
check} on the implementation and on floating-point idempotence of the
scan.
\end{block}

\begin{block}{Outputs}
\protect\phantomsection\label{outputs}
For each run, structured JSON records reference summaries, per-constant
distances, ranks, and optional \texttt{q\_sensitivity} blocks; CSV
mirrors the primary constant table. Histograms (uniform subsample,
pooled subsample, and \(\mu_Q\) on the pooled subsample) are written
under \texttt{output/figures/} for visual comparison of tail placement.
\end{block}
\end{frame}

\end{document}
